Lograr que un modelo de lenguaje convencional funcione como un tutor socrático que sepa limitarse depende de un diseño de prompts muy cuidado. No basta con darle un papel a la IA sino que hay que estructurar sus instrucciones para que use métodos como la dialéctica o la mayéutica \cite{chang2023prompting}. El principal obstáculo para conseguirlo es que los modelos actuales están diseñados por defecto para actuar como asistentes que ofrecen respuestas directas y rápidas a cualquier consulta. Esta naturaleza choca con el entorno educativo donde el objetivo es que el sistema actúe como una guía que no regale el código final. Para superar esta limitación las investigaciones proponen una evaluación jerárquica donde la IA primero identifica si la pregunta del alumno es general o específica y analiza el estado de su código antes de responder. Al aplicar esta lógica el modelo puede formular preguntas que parten del nivel de conocimiento real del estudiante para que sea él quien descubra la solución mediante su propio razonamiento \cite{kargupta2024instruct}.